\chapter{Реални числа}

\section*{Лекция 01 - 7 октомври 2026 г.}
За да можем да правим математически анализ, се нуждаем от знания за обекта, върху който работим: реалните числа. Това за съжаление е доста трудна задача. Строгото изграждане на тази теория и доказателството на двете теореми, с които ще започне следващата лекция, излизат извън този курс. Ще се задоволим със строга формулировка на основните свойства на реалните числа и интуитивно обяснение на тяхната природа.

\section{Основни понятия}

\begin{itemize}
    \item множества, означения, квантори -- ще ги оставим на курса по Дискретни структури;
    \item изображения (графика, инекция, сюрекция, биекция) -- оставяме ги на упражненията.
\end{itemize}

\section{Множеството от безкрайните десетични дроби}

Понятието за число се изгражда постепенно с развитието на човечеството като цяло и на всеки индивид в частност. От училище сте запознати с рационалните числа и можете да смятате с тях. Да фиксираме следните стандартни означения:
\begin{itemize}
    \item множеството на естествените числа $\mathbb{N} = \{1, 2, 3, 4, \dots\}$;
    \item множеството на целите числа $\mathbb{Z} = \{\dots, -3, -2, -1, 0, 1, 2, 3, \dots\}$;
    \item множеството на рационалните числа $\mathbb{Q} = \left\{ \dfrac{m}{n} : m \in \mathbb{Z}, n \in \mathbb{N} \right\}$.
\end{itemize}

Също от училище сте свикнали с подреждането на числата, с които се запознавате, върху „числовата права“. По този начин се получава съответствие между геометрични обекти (точки върху правата) и числа. Естествено е да очаквате, че трябва да можете да „мерите“ дължината на дадена отсечка, тоест да ѝ съпоставите някакво неотрицателно число по един разумен начин, съгласуван с опита. Още в древна Гърция е забелязано, че дължината на диагонала на квадрат „не е съизмерима“ с дължината на неговата страна, тоест не е рационално число, умножено по дължината на квадрата. Ако квадратът е със страна едно, от теоремата на Питагор знаем, че дължината на диагонала, повдигната на квадрат, трябва да е две.

\begin{problembox}{Пример 2.1}
Числото $\sqrt{2}$ (дължината на диагонала на квадрат със страна $1$) не е рационално, т.е. $\sqrt{2} \notin \mathbb{Q}$.
\end{problembox}

\begin{proof}
Допускаме противното, че $\sqrt{2} \in \mathbb{Q}$, по-точно $\sqrt{2} = \dfrac{p}{q}$ за $p, q \in \mathbb{N}$. Без ограничение на общността (б.о.о.) считаме, че $\operatorname{HOD}(p, q) = 1$, т.е. числата $p, q$ са взаимно прости. Тогава:
\[
\sqrt{2}q = p \implies (\sqrt{2}q)^2 = p^2 \implies 2q^2 = p^2.
\]
Понеже $2 \mid p^2$, то следва $2 \mid p$, или $p = 2r$ за някое $r \in \mathbb{N}$. Заместваме $p$:
\[
2q^2 = (2r)^2 = 4r^2 \implies q^2 = 2r^2.
\]
Оттук имаме $2 \mid q^2 \implies 2 \mid q$. Оказва се, че $p$ и $q$ са четни, което е противоречие с допускането, че $\operatorname{HOD}(p, q) = 1$.
\end{proof}

Примери като този демонстрират необходимостта от разширяване на понятието за число. Надявам се, че този курс ще покаже, че тази необходимост далеч не се изчерпва с възможността да намираме корени на полиноми с цели коефициенти (например $x^2 - 2 = 0$).

Като най-близък до интуицията, която имате от досегашния си опит, е моделът на реалните числа като множеството $\mathbb{R}$ от всички безкрайни десетични дроби. От физическа гледна точка, ако „мерим“ някаква физическа величина, това винаги става с някаква (все по-голяма с развитието на технологиите) точност. Нека мерим величината $a$ (дължина, тегло, време или нещо друго):
\begin{enumerate}
    \item Ако мерим с точност до $1$, то съществува $a_0 \in \mathbb{N} \cup \{0\}$, такова че: $a_0 \le a \le a_0 + 1$.
    \item Ако мерим с точност до $0.1$, то съществува $a_1 \in \{0, 1, 2, \dots, 9\}$, такова че $a_0.a_1 \le a \le a_0.a_1 + 0.1$.
    \item Ако мерим с точност до $0.01$, то съществува $a_2 \in \{0, 1, 2, \dots, 9\}$, такова че $a_0.a_1a_2 \le a \le a_0.a_1a_2 + 0.01$ и т.н.
\end{enumerate}

Тук идва абстракцията: мислим, че можем да продължаваме този процес до безкрайност и да представим величината $a$ като безкрайната десетична дроб $a_0.a_1a_2 \cdots a_n \cdots$.

И тъй, означаваме множеството от всички безкрайни десетични дроби с:
\[
\mathbb{R} = \{ s\, a_0.a_1a_2\cdots : s \in \{+, -\}, a_0 \in \mathbb{N} \cup \{0\}, a_n \in \{0, 1, \dots, 9\}, n \in \mathbb{N} \}.
\]

Горното означение не е съвсем точно: възможно е две различни безкрайни десетични дроби да представят едно и също число (такъв ефект има и при $\mathbb{Q}$: $1/2 = 3/6$). Първият пример е $0.999\cdots = 1.000\cdots$. Опитайте се да се убедите в това. Формално, две различни безкрайни десетични дроби $a, b \in \mathbb{R}$ представят едно и също число при следните условия:

\begin{definition}{Равенство на $a, b \in \mathbb{R}$}{def:equality_reals}
Нека $a = a_0.a_1a_2\cdots$ и $b = b_0.b_1b_2\cdots$. Казваме, че $a$ и $b$ са равни като реални числа, ако имат еднакви знаци (знаците могат да са различни само ако и двете дроби са нулеви) и:
\begin{itemize}
    \item или $a_i = b_i$ за всички $i \in \{0, 1, 2, \dots\}$ (тоест десетичните им представяния съвпадат);
    \item или съществува $n \in \mathbb{N} \cup \{0\}$ такова, че $a_i = b_i$ за всички $i \in \{0, 1, 2, \dots, n - 1\}$, $a_n = b_n + 1$ и $a_j = 0$, $b_j = 9$ за всички $j \in \{n + 1, n + 2, \dots\}$ (тоест $a_0.a_1a_2\cdots a_{n-1}a_n 000\cdots = a_0.a_1a_2\cdots a_{n-1}(a_n - 1)999\cdots$).
\end{itemize}
\end{definition}

От училище ви е известно, че всяко рационално число може да бъде представено като крайна или безкрайна периодична дроб, значи имаме $\mathbb{N} \subset \mathbb{Z} \subset \mathbb{Q} \subset \mathbb{R}$.

\begin{definition}{Сравняване на $a, b \in \mathbb{R}$}{def:comparison_reals}
\begin{enumerate}
    \item Ако $a \ge 0, b \ge 0$, то:
    \[
    a \le b \iff
    \begin{cases}
    a_i = b_i, \quad \forall i \in \{0, 1, 2, \dots\} \\
    \text{или} \\
    \exists n \in \mathbb{N} \cup \{0\} : a_i = b_i \ \forall i \in \{0, 1, \dots, n - 1\} \text{ и } a_n < b_n.
    \end{cases}
    \]
    \item Ако $a \le 0, b \ge 0$, то $a \le b$.
    \item Ако $a \le 0, b \le 0$, то $a \le b \iff |a| \ge |b|$.
\end{enumerate}
\end{definition}

Ако $a \le b$ и $a, b$ не са равни, то пишем $a < b$. По този начин вече можем да сравняваме безкрайни десетични дроби или, иначе казано, в множеството $\mathbb{R}$ има дефинирана наредба.

За да говорим за безкрайните десетични дроби като за числа, е важно и да можем да извършваме алгебричните операции събиране и умножение с тях (и при това основните закони за събирането и умножението да остават в сила). Да, това може да се направи, и тези, които се интересуват, могат да намерят съответните формални дефиниции и строга проверка на законите в учебника на Садовничий и Сендов. Доста е дълго, защото винаги може да има пренос, който да промени десетичния запис на началото. Неформалната идея обаче е ясна: „отрязваме“ безкрайните десетични дроби, които ще събираме (или умножаваме) до $n$-тия знак след десетичната точка и събираме (или умножаваме) получените крайни десетични дроби. Идеята е, че сумата на крайните десетични дроби е близо (колко близо зависи от $n$) до сумата на безкрайните десетични дроби.

\section{Ограничени множества от реални числа, горна и долна граница, супремум и инфимум}

Сега можем да забравим за природата на елементите на $\mathbb{R}$ и да работим само с тяхната наредба. Основните свойства на наредбата са описани в началото на следващата лекция.

Дефинициите от тази секция са първото ново знание, което ще бъде постоянно използвано през целия курс по ДИС1 (също и ДИС2). В тези дефиниции $A$ е подмножество на $\mathbb{R}$ (и е важна само наредбата в $\mathbb{R}$).

\begin{definition}{Множество, ограничено отгоре. Множество, ограничено отдолу}{def:bounded_sets}
Казваме, че $A$ е \textbf{ограничено отгоре}, ако съществува $M \in \mathbb{R}$ такова, че всички елементи $a$ на $A$ са по-малки или равни на $M$ ($\exists M \in \mathbb{R} \ \forall a \in A : a \le M$). Ако $M \ge a$ за всяко $a \in A$, то $M$ се нарича \textbf{горна граница (мажоранта)} на $A$. Друг начин да запишем същото е $A \subset (-\infty, M]$.

Съвсем аналогично се дефинира множество, \textbf{ограничено отдолу}: $A$ е ограничено отдолу, ако съществува $M \in \mathbb{R}$ такова, че всички елементи $a$ на $A$ са по-големи или равни на $M$ ($\exists M \in \mathbb{R} \ \forall a \in A : a \ge M$). Всяко $M \in \mathbb{R}$, за което $A \subset [M, +\infty)$, се нарича \textbf{долна граница (миноранта)} на $A$.
\end{definition}

Едно от нещата, на които ще се учим, е умението за отрицание (negation) на формални твърдения. В следващите няколко реда ще упражним именно него.

\begin{problembox}{Пример 3.2}
$A$ не е ограничено отгоре точно когато за всяко реално число съществува елемент на $A$, който е по-голям от числото:
\begin{align*}
\neg (\exists M \in \mathbb{R} \ \forall a \in A : a \le M) &\equiv \forall M \in \mathbb{R} \ \neg (\forall a \in A : a \le M) \\
&\equiv \forall M \in \mathbb{R} \ \exists a \in A : \neg (a \le M) \\
&\equiv \forall M \in \mathbb{R} \ \exists a \in A : a > M.
\end{align*}
\end{problembox}

\begin{definition}{Ограничено множество}{def:bounded_set}
$A$ се нарича \textbf{ограничено}, ако $A$ е ограничено отгоре и $A$ е ограничено отдолу.
\end{definition}

Следната дефиниция е централната за тази лекция:

\begin{definition}{Точна горна граница. Точна долна граница}{sup_inf}
Нека $A \subset \mathbb{R}$ е непразно множество, ограничено отгоре. \textbf{Точна горна граница} или \textbf{супремум} на $A$ ($\sup A$) наричаме най-малката от всички горни граници на $A$, т.е.:
\begin{enumerate}
    \item $a \le \sup A, \ \forall a \in A$ ($\sup A$ е горна граница на $A$).
    \item $\forall c < \sup A \ \exists a \in A : a > c$ (няма по-малка горна граница от $\sup A$).
\end{enumerate}

Нека $A \subset \mathbb{R}$ е непразно множество, ограничено отдолу. \textbf{Точна долна граница} или \textbf{инфимум} на $A$ ($\inf A$) наричаме най-голямата от всички долни граници на $A$.
\end{definition}

\begin{problembox}{Пример 3.5}
Да видим някои примери за супремуми на множества:
\begin{itemize}
    \item Супремумът на $[0, 1]$ е $\sup [0, 1] = 1$;
    \item $\sup (0, 1) = 1$;
    \item $\sup \left\{ 1 - \dfrac{1}{n} : n \in \mathbb{N} \right\} = 1$;
    \item $\sup \left( (-2, 0) \cup \left\{ \dfrac{1}{2} \right\} \right) = \dfrac{1}{2}$.
\end{itemize}
\end{problembox}

\begin{problembox}{Пример 3.6}
Какво означава, че $c$ не е $\sup A$? Това е вярно точно когато $\exists a \in A : a > c$ ($c$ не е горна граница за $A$) или $\exists d < c \ \forall a \in A : a \le d$ (съществува по-малка горна граница). Използвахме:
\[
\neg (\forall c < \sup A \ \exists a \in A : a > c) \equiv \exists c < \sup A \ \forall a \in A : a \le c.
\]
\end{problembox}

\begin{problembox}{Упражнение 3.7}
Докажете, че $\sup (A + B) = \sup A + \sup B$, където сме дефинирали:
\[
A + B := \{a + b : a \in A, b \in B\}.
\]
\end{problembox}

\begin{proof}
Трябва да проверим, че сборът от супремумите на $A$ и $B$ е горна граница на множеството $A + B$, от една страна, а от друга -- да видим, че всяко $c < \sup A + \sup B$ не се явява горна граница на $A + B$.

Първо ще се убедим, че $\sup A + \sup B$ е горна граница на $A + B$. Ползваме означенията $a \in A, b \in B$. Нека $x = a + b \in A + B$ е произволен елемент на множеството $A + B$. Понеже:
\[
\begin{cases}
a \in A, \sup A \ge a \\
b \in B, \sup B \ge b
\end{cases}
\implies x = a + b \le \sup A + \sup B.
\]
Следователно $\forall x \in A + B : \sup A + \sup B \ge x$ и значи $\sup A + \sup B$ е горна граница за $A + B$.

Сега взимаме произволно $c < \sup A + \sup B$. Да положим $\varepsilon = \sup A + \sup B - c > 0$. Тогава $\dfrac{\varepsilon}{2} = \dfrac{\sup A + \sup B - c}{2} > 0$ и следователно имаме:
\[
\begin{cases}
\sup A - \dfrac{\varepsilon}{2} < \sup A \implies \exists a_0 \in A, \ a_0 > \sup A - \dfrac{\varepsilon}{2} \\[2mm]
\sup B - \dfrac{\varepsilon}{2} < \sup B \implies \exists b_0 \in B, \ b_0 > \sup B - \dfrac{\varepsilon}{2}
\end{cases}
\]
Следователно $a_0 + b_0 > \sup A + \sup B - \varepsilon$. Понеже $\sup A + \sup B - \varepsilon = c$, показахме, че $a_0 + b_0 > c$, като $x_0 := a_0 + b_0 \in A + B$, и следователно $c$ няма как да бъде горна граница на $A + B$.
\end{proof}

Ще докажем, че в множеството от безкрайните десетични дроби е в сила следното изключително важно свойство:

\begin{theorem}{Принцип за непрекъснатост}{continuity_principle}
Всяко ограничено отгоре непразно множество $A \subset \mathbb{R}$ има супремум. Всяко ограничено отдолу непразно множество $A \subset \mathbb{R}$ има инфимум.
\end{theorem}

\begin{problembox}{Пример 3.9}
Ако $A = \{ q \in \mathbb{Q} : q \ge 0, q^2 \le 2 \}$, то $\sup A = \sqrt{2} \notin \mathbb{Q}$.
\end{problembox}

\begin{definition}{Цяла част на $x \in \mathbb{R}$}{floor}
Нека $x \in \mathbb{R}$ е произволно. \textbf{Цяла част} на числото $x$ наричаме най-голямото цяло число $z \in \mathbb{Z}$, ненадвишаващо $x$. Бележим го с $\lfloor x \rfloor$. Формално:
\[
\lfloor x \rfloor = \max \{z \in \mathbb{Z} : z \le x\}.
\]
\end{definition}

За тази функция ($\lfloor \cdot \rfloor : \mathbb{R} \to \mathbb{Z}$) могат да се съставят много примери и лесно да се начертае характерната стъпаловидна графика. Така например, $\lfloor -\pi \rfloor = -4$; $\lfloor \pi \rfloor = 3$; $\lfloor 0 \rfloor = 0$.
\drawing{../figures/floor_function.pdf}{Графика на функцията цяла част $y = \lfloor x \rfloor$}

\begin{proof}[Доказателство на Твърдение \ref{thm:continuity_principle}]
Първо да разгледаме случая, когато $B := A \cap [0, +\infty)$ не е празно. Ясно е, че супремумите на $A$ и $B$, ако съществуват, трябва да съвпадат. Ясно е също, че $B$ е ограничено отгоре. Тогава $\{\lfloor a \rfloor : a \in B\}$ е ограничено отгоре подмножество на $\{0\} \cup \mathbb{N}$, следователно е крайно.

\begin{enumerate}
    \item Избираме $\overline{a}_0 = \max \{\lfloor a \rfloor : a \in B\}$ -- най-голямата цяла част на число, което е елемент на $B$.
    
    Нека сега означим $B_1 = \{a \in B : a = \overline{a}_0.a_1a_2\cdots\} \neq \emptyset$ -- реалните числа от $B$ с цяла част точно $\overline{a}_0$.

    \item Избираме $\overline{a}_1 = \max \{a_1 \in \{0, 1, \dots, 9\} : \exists a \in B_1, a = \overline{a}_0.a_1a_2\cdots\}$ -- това е максималното такова число измежду $0, 1, \dots, 9$, за което има елемент от $B$ от горния вид. В някакъв смисъл -- най-голям елемент на множеството, ако сравняваме елементите само по $a_0$ и $a_1$ (всички елементи са от вида $a = a_0.a_1a_2\cdots$, като в записа участва десетична точка).

    Аналогично строим $B_2 = \{a \in B_1 : a = \overline{a}_0.\overline{a}_1a_2\cdots\} \neq \emptyset$.

    \item Продължаваме с $\overline{a}_2 = \max \{a_2 \in \{0, 1, \dots, 9\} : \exists a \in B_2, a = \overline{a}_0.\overline{a}_1a_2\cdots\}$ и т.н.
\end{enumerate}

Така можем да построим $\overline{a} = \overline{a}_0.\overline{a}_1\overline{a}_2\cdots \in \mathbb{R}$. Твърдим, че $\overline{a} = \sup B$.

Първо да проверим, че $\overline{a} \ge a$ за всяко $a \in B$. Фиксираме произволно $a \in B$. Ако $a_i = \overline{a}_i$ за всяко $i \in \{0, 1, 2, \dots\}$, то $\overline{a} = a$. В противен случай съществува $n \in \mathbb{N} \cup \{0\}$ такова, че $a_i = \overline{a}_i$ за всички $i \in \{0, 1, \dots, n - 1\}$ и $a_n \neq \overline{a}_n$. Случаят $a_n > \overline{a}_n$ е невъзможен от начина на построение на $\overline{a}_n$. Следователно $a_n < \overline{a}_n$, откъдето имаме $\overline{a} \ge a$. С това проверихме, че $\overline{a}$ е горна граница за $B$.

Нека сега $d < \overline{a}$ е произволно. Необходимо е да покажем, че $d < a$ за някое $a \in B$, позовавайки се на Дефиниция \ref{def:sup_inf}.

Тъй като $d < \overline{a}$, съществува $n \in \mathbb{N}$ такова, че:
\[
\begin{cases}
d_0 = \overline{a}_0 \\
d_1 = \overline{a}_1 \\
\quad \dots \\
d_n = \overline{a}_n \\
d_{n+1} < \overline{a}_{n+1}.
\end{cases}
\]
Но $\overline{a}_{n+1} = \max \{a_{n+1} \in \{0, 1, \dots, 9\} : \exists a \in B_n, a = \overline{a}_0.\overline{a}_1\overline{a}_2\cdots\overline{a}_n a_{n+1}\cdots\}$. Следователно съществува $a \in B$, за което $a = \overline{a}_0.\overline{a}_1\overline{a}_2\cdots\overline{a}_n \overline{a}_{n+1} a_{n+2}\cdots$ и значи $a > d$. Следователно $d$ не е горна граница за $B$.

Ако $A \subset (-\infty, 0)$, прилагаме подобна конструкция към $\{|a| : a \in A\}$, но вместо максимални вземаме минимални елементи и полученото число вземаме със знак минус.
\end{proof}

\subsection*{Задачи за обмисляне вкъщи (до следващата седмица)}

\begin{problembox}{Упражнение 3.11}
Да се докаже, че $A$ е ограничено точно тогава, когато съществува $M \ge 0$ такова, че $|a| \le M$ за всяко $a \in A$.
\end{problembox}

\begin{problembox}{Упражнение 3.12 (Принцип на трансфинитната индукция)}
Нека $a \in \mathbb{R}$ и нека $M \subset [a, +\infty)$ е непразно множество от реални числа. Нека за всяко $b \ge a$ знаем, че ако $[a, b) \subset M$, то съществува $\varepsilon > 0$ такова, че $[a, b + \varepsilon) \subset M$. Да се докаже, че $M = [a, +\infty)$.

\textit{Забележка:} В горното при $b = a$ формално получаваме $[a, a) = \emptyset \subset M$, което влече $[a, a + \varepsilon) \subset M$ за някое положително $\varepsilon$.
\end{problembox}